% WHAT AN ERROR SHOWS OF A MACRO IT IS READING.
%
% A fault raised while a macro's body is being read draws the macro in the
% context: its name, the text it matched, `->', and its body, with the
% reading marked where it stands. A body that asks for an argument draws a
% frame of its own above that one, `<argument>', holding what the argument
% was.
%
% trip.tex line 329 defines a macro with a delimited parameter,
%
%   \def\d#1\d{#1#1}
%
% and line 333 puts `\d#\d' where an alignment preamble is being read. The
% preamble scanner does not expand, so \d is not expanded there: it is
% expanded by the LOOKAHEAD OF THE GLUE SCAN in front of it -- `\tabskip1ex
% plus7200bp minus 4\wd4' has read a dimension and is looking to see whether
% another unit follows. The expansion puts `##' in the preamble, which is
% the second `#' of that tab, and the error is raised while the body is
% being read.
%
% The reference draws two frames, and so does this engine now:
%
%   ! Only one # is allowed per tab.
%   <argument> ##
%
%   \d #1\d ->#1#1
%
% It used to draw one, `\d \d ->####' -- both halves the substituted text
% rather than the definition, since `####' is what `#1#1' comes to once the
% argument `#' has been put in twice and each drawn doubled.
%
% WHY IT WAS HARD. The two engines expand a macro differently. The reference
% pushes the body AS IT STANDS and walks it; when it reaches `#1' it pushes
% the argument as a frame of its own, which is what `<argument>' names. This
% engine builds the substituted text first and pushes that, so by the time
% anything is read there is no body to show and no argument frame to show
% above it.
%
% WHAT WAS DONE, and what it costs: nothing is copied that was not copied
% before. The expansion is built in the input stack's store as it always
% was, with ONE WORD PER PARAMETER in front of it saying how long that
% argument came out, and the frame holds the definition. That is enough for
% both frames: walking the body against those lengths says how much of the
% BODY the reading has covered and which argument it stands in, and that
% argument's own tokens are inside the substituted text already, so they
% need not be kept anywhere else. The words are given back with the frame.
%
% The other way -- reading the body where it stands and pushing a frame per
% `#n', as the reference does -- was costed and not taken; the measurements
% are below.
%
% WHAT EACH WOULD COST, MEASURED. The counters `body parameter occurrences
% per call' and `body length per call' were added to settle this. They have
% now been read and taken out again; what they said, over the corpus rather
% than over trip, is this. testmath -- 41 pages of amsmath, the heaviest
% document in tests/corpus -- makes 21,204,416 macro calls, of which
% 9,708,986 substitute nothing. So 11,495,430 calls, 54%, are the ones the
% choice is about, and between them they ask for AT LEAST 31.7 million
% parameter occurrences (counting the `11 or more' bucket as eleven).
%
% Where the time goes, from a SIGPROF sampler over that same run, 17,652
% samples in the engine, -O2 with symbols. instantiate_macro is 25.0% of the
% whole run -- the largest single cost in the engine -- and divides:
%
%   5.67%  the counters themselves (temporary, now gone)
%   3.60%  matching the arguments        both options pay this
%   8.29%  the no-substitution fast path both options keep this
%   1.10%  working out the length        only the substituting one pays
%   5.61%  BUILDING THE SUBSTITUTION     only the substituting one pays
%   0.72%  teardown
%
% So reading the body where it stands can save at most about 6.7% -- the
% length and the copy -- and pays for it with 31.7 million frame pushes and
% pops in place of 11.5 million expansions built. The frame machinery is
% already about 5.7% of the run at its present volume
% (hstex_source_push_one, pop_frame, pop_exhausted_token_frames), so tripling
% the traffic through it plausibly costs more than the 6.7% it buys. That is
% an argument against the reference's own arrangement here, and it rests on
% this engine copying tokens cheaply and pushing frames dearly.
%
% Keeping the arguments alive and rebuilding the two frames only when a
% context is drawn costs nothing on the path that reads tokens, which is
% where the 25% is. It costs holding each call's arguments until its frame
% pops rather than until the call returns -- and frames pop in the order
% they were pushed, so the arena stays a stack and the extra it holds is
% bounded by how deep macros nest.
%
% On that evidence the second way is the one to take, and the first should
% not be attempted without measuring the frame traffic it would add.
%
% WHAT THE COUNTERS COST. They scanned every body on every call, which is
% why they are gone: interleaved A/B, minimum of user+sys over seven rounds
% on testmath, 17.07s with them and 16.41s without, so 3.87% of the run. The
% sampler put them at 5.67%; both agree they cost more than several of the
% things they were measuring.
%
% WHAT IS DRAWN HERE. [A] is the trip case cut down to its own file: the
% same definition, the same glue lookahead, the same error. [B] is the same
% macro under \tracingmacros, which hstex already draws exactly as the
% reference does -- so it is the context that is wrong and not the drawing
% of a macro. [C] takes two arguments undelimited and puts the second in
% twice, so that the reading stands inside the second of them when the fault
% is raised. The frame is named plainly `<argument>', not by number, and
% holds that argument's tokens; both engines now draw
%
%   <argument> ##
%
%   \e #1#2->#1#2#2
%
% where this engine used to draw `\e ->q####' -- `#1#2#2' with `q' put in
% for the first and `#', drawn doubled, put in twice for the second.
%
% [C] is also what says the argument found must be the one the reading is
% STANDING IN rather than the last one substituted: `#1' is `q' and `#2' is
% `#', and it is `#2' the frame must name.
\catcode`\{=1 \catcode`\}=2 \catcode`\#=6
\tracingonline=1 \errorcontextlines=100
\message{[A the trip case: an error while a substituted body is read]}
\def\d#1\d{#1#1}
\halign{\tabskip1pt\d#\d\cr}
\message{[B the same macro as tracingmacros draws it, which already agrees]}
\tracingmacros=2
\setbox0=\hbox{\d A\d}
\tracingmacros=0
\message{[C which argument the context names]}
\def\e#1#2{#1#2#2}
\halign{\tabskip1pt\e q#\cr}
\message{[done]}
\end
